(General background knowledge) Tato část popisuje praktické nasazení autograderů a kroky pro zavedení v kurzu.
(General background knowledge) Účel nasazení: zrychlit rutinní hodnocení, poskytovat okamžitou zpětnou vazbu a uvolnit čas vyučujícímu.
(General background knowledge) Základní kroky zavedení (krátký postup):
\begin{enumerate}
\item (General background knowledge) Definujte vhodné typy úloh a připravte testovací sady včetně okrajových případů; dělejte peer review testů.
\item (General background knowledge) Vyberte nástroj kompatibilní s LMS, integrujte ho do workflow odevzdání a pilotujte na malé skupině, sbírejte chyby a feedback.
\item (General background knowledge) Nastavte pravidlo lidského přezkumu pro sporné či otevřené odpovědi a iterujte podle zjištěných chyb.
\end{enumerate}
(General background knowledge) Měřte přínos (čas ručního hodnocení, opravy testů, spokojenost studentů, A/B piloty); mitigujte rizika (peer review testů, povinný lidský přezkum, testovací prostředí a IT podpora) a začínejte malými piloty.